% Generated draft: requires mathematical and source review.
\documentclass[11pt]{article}
\usepackage[margin=25mm]{geometry}
\usepackage{fontspec}
\setmainfont{Latin Modern Roman}
\usepackage{amsmath,amssymb,mathtools,unicode-math}
\setmathfont{Latin Modern Math}
\usepackage{enumitem,needspace,xurl,hyperref}
\setlength{\parindent}{0pt}
\setlength{\parskip}{5pt}
\setlength{\emergencystretch}{3em}
\newcommand{\N}{\mathbb{N}}
\newcommand{\Z}{\mathbb{Z}}
\newcommand{\Q}{\mathbb{Q}}
\newcommand{\R}{\mathbb{R}}
\newcommand{\C}{\mathbb{C}}
\newcommand{\F}{\mathbb{F}}
\newcommand{\A}{\mathbb{A}}
\newcommand{\PP}{\mathbb{P}}
\newcommand{\E}{\mathbb{E}}
\newcommand{\eps}{\varepsilon}
\newcommand{\dd}{\,\mathrm{d}}
\begin{document}
\section*{UnsolvedMath ID: 10400048}
\textbf{Source problem code:} catalogue entry
\subsection{Definitions and mathematical statement}
Let \(U\subset S^3\) be the unknot. The source asks to describe, for each degree \(d\), an abelian group represented by
\[
 \mathcal A_d=
 \{(M,K):(M,K)\sim_{HL_d}(S^3,U)\}\big/\sim_{HL_{d+1}}.
\]
This notation retains the extracted label \(HL_d\); the exact typography and definition of the source's relation must be checked.

More explicitly, after the admissible pair category \(\mathcal P\) and relation are fixed, the numerator is the \(HL_d\)-class of the base pair inside \(\mathcal P\), and the denominator identifies its elements by the next-degree relation. Determine the resulting abelian group, for example by an explicit isomorphism description or generators and relations. No finite generation or proposed rank is supplied or assumed.

The excerpt does not define \(HL_d\)-equivalence, its generating operations, degree range, or refinement from \(d\) to \(d+1\). It also does not state the conditions on \(M\), knot orientation/framing, or the binary operation making the quotient a group. Connected sum is a natural possibility for knot-pair constructions, but is not silently declared to be the intended operation. The abelian-group description is therefore a source assertion awaiting those definitions, not a structure derived here.

The supplied summary labels the question open/triage and says it was not obviously resolved, with no exact theorem or modern citation. This qualified assessment is retained without independent current verification.
\subsection{Short English statement}
Describe the next-degree quotient of knot-manifold pairs that are \(HL_d\)-equivalent to the unknot in \(S^3\). The equivalence relation, admissible pairs, degree range and group operation are missing from the excerpt.
\subsection{Sources}
\begin{itemize}
\item Local registry: \url{data/unsolvedmath/problems.json}, permanent UnsolvedMath ID \(10400048\), source code catalogue entry; supplied statement, background, and research summary, updated 2026-08-06.
\item Ohtsuki, \emph{Problems on invariants of knots and 3-manifolds} (2002), Problem 2.26, attributed to K. Habiro, PDF page 49: \url{https://msp.org/gtm/2002/04/gtm-2002-04-024s.pdf}. Supplied PDF-font-structure extraction accessed 2026-07-29; provenance distinguishes 231 numbered items.
\item No separate relation definition, group presentation or resolving publication is supplied.
\end{itemize}
\end{document}
